\scanpage{140}
\[
\begin{aligned}
\mathbb N\deleted{_{\unsure{x}}}&\longrightarrow I: \\[-2pt]
&\phantom{\longrightarrow}\{i\}
\end{aligned}
\]
\[
 I\longrightarrow\unsure{\mathscr D|\mathscr O}.
\]
\[
\left.\begin{aligned}
 &\unsure{\mathbb I|I}\ \text{filtrante finie}\\
 &\mathbb N\longrightarrow I
\end{aligned}\right\}\quad
\begin{aligned}
 &\text{alors $I$ est grand devant $\pi$}\\
 &\text{pour tt cardinal $\pi$}
\end{aligned}
\]
\deleted{\unsure{car pour tt} \illegible{1}}

\uline{Prop} Soit $I$ une catégorie équiv. à une petite cat. Conditions équivalentes, soit $\pi_0=\operatorname{card}\operatorname{Fl} I$

a) $I$ est grand devant $\pi$ pour tt card $\pi$

b) $I$ est grand devant $\pi'_0=\operatorname{Sup}(\pi_0,\aleph_0)$ (i.e. \unsure{$I$ grand} devant $\pi_0$, si $\pi_0$ est infini)

c) Il existe un foncteur cofinal $I'\longrightarrow\mathbb I$, où $I'$ est une catégorie \unsure{filtrante} finie.

d) \deleted{Il existe un foncteur \unsure{cofinal} $I'$} \inserted{comme c), \unsure{sauf|avec|mais} $I'$ \deleted{\unsure{étant}} la cat.} $I^{\unsure{*}}_0$ réduite à un seul objet, avec une seule flèche non identique $\unsure{p|\rho}$, satisfaisant $p^2=p$.

e) Il existe un objet $x_0$ de $I$, et une section de $x_0$ sur l'objet final $e_{I^\wedge}$ de $I^\wedge$ -- i.e. pour tt $x\in\operatorname{Ob} I$, un $x\xrightarrow{f_x}x_0$ tel que .... \unsure{sinon pas}.

\deleted{\unsure{réciproquement} \illegible{2}}
\diagram{\begin{tikzcd}[ampersand replacement=\&,row sep=small,column sep=small]
\deleted{a}\arrow[r,Rightarrow]\&\deleted{\unsure{b}}\arrow[r,Rightarrow]\&\deleted{\unsure{c}}\arrow[r,Rightarrow]\&\deleted{\unsure{d}}\\
\deleted{c}\arrow[r,Rightarrow]\&\deleted{\unsure{a}}\arrow[r,Rightarrow]\&\deleted{b}\&
\end{tikzcd}}
\margin{\deleted{\unsure{démonstration}} pour une cat. finie}
\diagram{\begin{tikzcd}[ampersand replacement=\&,row sep=small,column sep=small]
\&\&\&\&c\\
e\arrow[r,Leftrightarrow]\&d\arrow[r,Rightarrow]\&a\arrow[r,Rightarrow]\&b\arrow[ur,Rightarrow]\arrow[r,Rightarrow]\&e
\end{tikzcd}}
\margin{$\downarrow$ \unsure{exc.|exerc.} \unsure{dét.|déterm.} des foncteurs cofinaux $I_0\to I$}
\margin{$\searrow$ \deleted{\unsure{tr}} \unsure{triv.}}
\[
\begin{aligned}
 &\unsure{\text{donc}}\ e\Longleftrightarrow d\Longleftrightarrow a\Longrightarrow b\\
 &\deleted{\unsure{b\Longrightarrow}}\Longrightarrow c
\end{aligned}
\]
\margin{\deleted{\unsure{on}} \unsure{Appliquant|Appliquer} ceci au cas d'une catégorie \deleted{finie} $I'$, on trouve \deleted{\unsure{que}} \deleted{\unsure{b}} $c\Rightarrow d$ dans tt cat équiv. \unsure{une}}

\diagram{\begin{tikzcd}[ampersand replacement=\&,row sep=small,column sep=small]
{}\arrow[r,shift left=1.2ex]\arrow[r,shift left=.4ex]\arrow[r,shift right=.4ex]\arrow[r,shift right=1.2ex]\&x\arrow[r,"p"]\&\deleted{\unsure{x}}
\end{tikzcd}}
\[
\begin{aligned}pu&=u\\p^2&\deleted{\unsure{\ne}}=p\end{aligned}
\]
% margin note written upside down
\diagram{\begin{tikzcd}[ampersand replacement=\&,row sep=small,column sep=small]
\&\underline S(\unsure{\pi|\mathbb I})_{/x}\arrow[d,"\text{cofinal}"]\\
J\arrow[ur,bend left=14,"{\deleted{\text{4.4.2}}}"]\arrow[r,"{\deleted{\text{cof}}}","{\substack{\text{cof.}\\\text{par }4.4.2.}}"']\&\mathcal M_{\pi/x}
\end{tikzcd}}
% archive sheet number: 140
